\section{Clôture \(\beta\) de la section électronique et expression dimensionnelle de son énergie au repos}
\label{chap:biografia-electron-rev6}
\index{électron!construction holographique}
\index{positron!bande conjuguée}
\index{masse!loi électronique}

L'électron apparaît dans HMT comme le point où s'assemblent six
constructions antérieures et mutuellement typées : l'incompatibilité entre
les deux feuillets arithmétiques d'APP, l'orientation induite par l'involution
de la charnière surfacique--volumétrique, deux quotients entiers du recensement APP--TPK,
la clôture
dodécaphasique de \(\alpha_{\mathrm{HMT}}\), le défaut global
\(\Delta_4\) et le centre unique de l'atlas nonadique. La bande de route ajoute
ensuite la conjugaison électron--positron. La carte en MeV et la comparaison
avec CODATA sont postérieures à toutes ces constructions.

Cette antériorité fait partie du contenu mathématique : les facteurs, leurs
types et leur ordre causal sont déterminés avant l'ouverture de la carte tabulée, qui
n'intervient que dans la comparaison métrologique ultérieure.

\subsection{Scalaire électronique \(\mathfrak e_{\mathrm{HMT}}\) et graphe causal de ses facteurs APP–TPK, \(\alpha_{\mathrm{HMT}}\), \(\Delta_4\) et atlas nonadique}

La sortie interne est le scalaire adimensionnel
\begin{equation}
 \label{eq:electron-holografico-rev6}
 \boxed{
 \mathfrak e_{\mathrm{HMT}}
 =
 \frac{\sqrt3}{4}
 \left[
 \exp\!\left(\frac{\varphi}{\pi^2}\right)
 -\frac{396}{18}\,\alpha_{\mathrm{HMT}}^3
 \right]
 \left[
 1+\frac{270}{18}\,\Delta_4
 \right].}
\end{equation}
Les identités
\[
 \frac{396}{18}=22=27-5,
 \qquad
 \frac{270}{18}=15=\binom62=3\cdot5
 \]
permettent de retrouver l'écriture abrégée
\[
 \mathfrak e_{\mathrm{HMT}}
 =
 \frac{\sqrt3}{4}
 \left(e^{\varphi/\pi^2}-22\alpha_{\mathrm{HMT}}^3\right)
 (1+15\Delta_4).
\]
Ici, la lettre \(e\) à l'intérieur de l'exponentielle est le nombre d'Euler ;
l'indice de \(\mathfrak e_{\mathrm{HMT}}\) désigne la section électronique.

Le graphe causal de l'équation est
\[
\begin{array}{ccl}
\text{projecteurs APP en dimension trois}
&\longrightarrow& \sqrt3/4,\\[2mm]
\text{calendrier TPK et feuillets aréal--volumétrique}
&\longrightarrow&
F_{\mathrm{av}},\ \pi/\varphi^2,\
\varphi/\pi^2,\\[2mm]
\text{clôture dodécaphasique}
&\longrightarrow&\alpha_{\mathrm{HMT}},\\[2mm]
\text{recensement APP--TPK et famille électronique}
&\longrightarrow&396,\ 18,\ 270,\ 22,\ 15,\ 5,\\[2mm]
\text{lectures globales de }\pi,e\text{ et registre}
&\longrightarrow&\Delta_4,\\[2mm]
\text{atlas nonadique et anti-section}
&\longrightarrow&(5,5),\
\text{bande électron--positron}.
\end{array}
\]
Aucune flèche de ce diagramme ne reçoit comme argument la masse électronique
de référence.

\subsection{1ᵉʳ facteur : incompatibilité exacte des feuillets APP}

Pour \(d\geq2\), soit
\[
 \mathscr A_d=\{1,\ldots,9\}^d
\]
muni de la mesure uniforme, et soient
\[
 \Sigma_d(x)=\operatorname{dr}(x_1+\cdots+x_d),
 \qquad
 \Pi_d(x)=\operatorname{dr}(x_1\cdots x_d)
\]
les observables additive et multiplicative. Leurs espérances conditionnelles
orthogonales sur \(\ell^2(\mathscr A_d)\) sont
\(P_{\Sigma,d}\) et \(P_{\Pi,d}\). La
\cref{prop:app-no-conmutatividad} démontre par énumération exacte que
\[
 \left\|[P_{\Sigma,3},P_{\Pi,3}]\right\|
 =\frac{\sqrt3}{4}.
\]

\begin{proposition}[Contenu APP du préfacteur électronique]
\label{prop:prefactor-app-electron-rev6}
Le préfacteur \(\sqrt3/4\) de
\cref{eq:electron-holografico-rev6} est la norme exacte de la
non-commutativité entre les feuillets somme et produit en dimension trois.
Ce n'est pas un facteur géométrique introduit à partir de la formule électronique.
\end{proposition}

\begin{proof}
Le tableau conjoint de \(\Sigma_3\) et \(\Pi_3\) détermine les angles
principaux entre
\(\operatorname{Ran}P_{\Sigma,3}\) et
\(\operatorname{Ran}P_{\Pi,3}\). Si \(s\) est l'une de ses valeurs
singulières, la contribution correspondante à la norme du commutateur est
\(s\sqrt{1-s^2}\). L'énumération des \(9^3\) mots donne pour maximum
\[
 s^2(1-s^2)=\frac3{16},
 \qquad s^2=\frac14.
\]
La norme maximale est donc \(\sqrt{3/16}=\sqrt3/4\), avant
d'introduire \(\alpha_{\mathrm{HMT}}\), \(\Delta_4\), une unité d'énergie
ou une masse de référence.
\end{proof}

La fonction physique du facteur est ainsi fixée : il mesure l'incompatibilité
maximale des deux lectures arithmétiques sur le même support ternaire.
L'électron n'hérite pas de deux nombres indépendants, l'un additif et l'autre
multiplicatif ; il hérite de l'obstruction exacte à rendre simultanément
compatibles les deux feuillets. Le préfacteur conserve dans la loi de masse la
mémoire locale du désaccord APP dont naît la route.

\subsection{2ᵉ facteur : l'involution d'échange entre \(\pi\) et \(\varphi\)}

Le calendrier primitif des modes
\((\Sigma,\Pi,\Pi)\) induit, dans les feuillets surfacique et volumétrique, la matrice
\[
 F_{\mathrm{av}}
 =
 \begin{pmatrix}
 0&1\\
 1&1
 \end{pmatrix}.
\]
Le \cref{thm:descenso-necesario-fav} l'obtient en comptant les transitions
\(\Sigma\to\Pi\), \(\Pi\to\Pi\) et \(\Pi\to\Sigma\), chacune une fois.
La matrice n'est pas postulée à partir du nombre d'or : ses valeurs propres
\(\varphi\) et \(-\varphi^{-1}\) apparaissent après la construction de
l'incidence.

Dans la mémoire de deuxième ordre,
\(\operatorname{Sym}^2(F_{\mathrm{av}})\) a pour spectre
\[
 \varphi^2,\qquad -1,\qquad\varphi^{-2}.
\]
Ainsi, \(\varphi^{-2}\) est le mode stable positif de la charnière. La marque
de clôture \(\pi\), appliquée une seule fois au retour surfacique, produit
\[
 \frac{\pi}{\varphi^2}.
\]
Ce rapport conserve l'ordre clôture--auto-échelle. La loi électronique utilise son
orientation opposée.

\begin{theorem}[Involution exacte de la charnière électronique]
\label{thm:espejo-electron-rev6}
Soient
\[
 \mathscr R(x,y)=\frac{x}{y^2},
 \qquad
 \mathsf J(x,y)=(y,x).
\]
Alors \(\mathsf J^2=I\) et
\[
 \mathscr R(\pi,\varphi)=\frac{\pi}{\varphi^2},
 \qquad
 (\mathscr R\circ\mathsf J)(\pi,\varphi)
 =\frac{\varphi}{\pi^2}.
\]
Par conséquent,
\[
 \frac{\varphi}{\pi^2}
 =
 \frac{1}{
 \varphi^3\bigl(\pi/\varphi^2\bigr)^2},
 \qquad
 e^{\varphi/\pi^2}
 =
 \exp\!\left[
 \frac{1}{
 \varphi^3\bigl(\pi/\varphi^2\bigr)^2}
 \right].
\]
\end{theorem}

\begin{proof}
Le premier couple d'identités s'obtient en appliquant
\(\mathscr R\) avant et après l'involution \(\mathsf J\). Pour le
reparamétrage,
\[
 \varphi^3
 \left(\frac{\pi}{\varphi^2}\right)^2
 =\frac{\pi^2}{\varphi},
\]
et son inverse est \(\varphi/\pi^2\). La dernière égalité découle de l'application de
l'exponentielle.
\end{proof}

La présence simultanée de \(\pi\) et \(\varphi\) possède donc une
signification structurelle précise. \(\pi/\varphi^2\) lit la clôture surfacique
sur la mémoire volumétrique quadratique ; \(\varphi/\pi^2\) lit
l'auto-croissance intérieure sur la clôture quadratique. Ce ne sont pas deux quotients
semblables choisis par commodité. Ce sont les deux orientations d'un même
lecteur, échangées par une involution exacte. Le terme
\(e^{\varphi/\pi^2}\) publie dans la section électronique l'orientation
intérieure de cette charnière.

\subsection{3ᵉ et 4ᵉ facteurs : les entiers 22, 15 et le centre 5}

Le recensement de la carte APP employée par la famille électronique fournit
trois sommes entières :
\[
 \Sigma(\mathrm{APP})=396,
 \qquad
 \Sigma(\mathrm{centro}\ 2\times2)=18,
 \qquad
 \Sigma(\mathrm{canal})=270.
\]
La première admet la décomposition en anneaux
\[
 396=18+72+108+198.
\]
Le noyau \(2\times2\) est donc une sous-structure explicite de la
carte complète, et l'entier 270 appartient au canal global du même recensement
APP--TPK.
La somme des anneaux pairs est
\[
 18+108=126,
 \qquad
 \frac{396}{126}=\frac{22}{7}.
\]
Ce dernier quotient est le sceau rationnel de clôture de la carte, non une
identification de \(22/7\) au nombre \(\pi\) engendré par le lecteur
coinductif. Il montre bien que le numérateur \(22\) remplit déjà une fonction
interne avant d'apparaître dans la loi électronique.

\begin{theorem}[Provenance interne des coefficients électroniques]
\label{thm:coeficientes-electron-rev6}
Dans la famille électronique fixée par la carte APP et son canal global,
\[
 \boxed{
 22=\frac{\Sigma(\mathrm{APP})}
 {\Sigma(\mathrm{centro})}
 =\frac{396}{18}=27-5,}
\]
\[
 \boxed{
 15=\frac{\Sigma(\mathrm{canal})}
 {\Sigma(\mathrm{centro})}
 =\frac{270}{18}=3\cdot5=\binom62.}
\]
En particulier, les coefficients sont déterminés sans consulter
\(Z_0\), la masse de l'électron ni une carte d'unités.
\end{theorem}

\begin{proof}
Les deux premières égalités sont des divisions exactes du recensement fini :
\[
 396=22\cdot18,
 \qquad
 270=15\cdot18.
\]
En définissant
\[
 p:=27-22,
\]
on obtient \(p=5\), et alors
\[
 15=3p=3\cdot5=\binom62.
\]
Toutes les quantités appartiennent au domaine entier de la carte et de la
famille avant la transduction métrologique.
\end{proof}

\subsubsection{Sélecteur de persistance dans le retour central}

Le même couple possède une construction dynamique qui fixe en outre son signe sans
consulter une masse. Les quatre orientations du germe central
\((5,5)\) reviennent pour la première fois à cette cellule après vingt-sept pas.
L'espace typé du retour est
\[
 \mathcal K_e=\mathbb Z/9\mathbb Z\times\mathbb F_3,
 \qquad |\mathcal K_e|=27.
\]
Par ailleurs, la microfibre de \(\pi\) contient les cinq régions
généalogiquement distinctes
\[
 \mathscr R_\pi=3_{++}+1_{--}+1_{\perp}.
\]
La jauge régionale \(\operatorname{diag}_c\) et, au sein d'une phase répétée,
l'ordre interne de \(U_6\) définissent l'injection
\[
 \iota:\mathscr R_\pi\hookrightarrow\mathcal K_e
\]
au moyen des cinq places distinctes
\[
 (4,0),\quad(5,0),\quad(6,0),\quad(6,1),\quad(6,2).
\]
Les deux premières correspondent respectivement à la région transversale et
au retour muté ; les trois dernières séparent les régions positives par leur
feuillet tritique. Le mot visible commun n'intervient pas comme inverse de l'application.

\begin{theorem}[Sélecteur interne des coefficients électroniques]
\label{thm:selector-persistencia-electron-rev7}
Une fois fixés le retour central et les cinq régions typées, les coefficients de
l'équation électronique sont
\[
 \boxed{
 n_e=-\bigl|\mathcal K_e\setminus\iota(\mathscr R_\pi)\bigr|=-22,
 \qquad
 m_e^{\rm tor}=\bigl|\mathscr R_\pi\times\mathbb F_3\bigr|=+15.}
\]
Le premier signe enregistre le déficit de frontière ; le second, la mémoire
torsionnelle retenue.
\end{theorem}

\begin{proof}
Les cinq places précédentes sont distinctes, de sorte que \(\iota\) est
injective et
\[
 |\mathcal K_e\setminus\iota(\mathscr R_\pi)|=27-5=22.
\]
Chacune des cinq régions conserve trois feuillets tritiques, donc
\[
 |\mathscr R_\pi\times\mathbb F_3|=5\cdot3=15.
\]
La règle d'orientation attribue un signe négatif aux états retirés de la
coquille et un signe positif à la mémoire retenue. La construction n'ouvre que le
centre, le retour de vingt-sept états et les identités régionales ; elle
n'ouvre ni une masse observée, ni une impédance, ni une signature étalonnée.
\end{proof}

Le nombre \(22\) mesure combien de noyaux centraux de poids \(18\) recompose la
somme APP \(396\). Le nombre \(15\) mesure combien de ces mêmes noyaux
recompose le canal \(270\), et compte simultanément les paires d'une
hexade. Le nombre \(5\) n'est pas un paramètre supplémentaire :
\[
 5=27-22,\qquad 15=3\cdot5.
\]
Les trois écritures montrent la même clôture depuis la profondeur 27,
depuis la normalisation centrale et depuis la combinatoire hexadique.

\subsubsection{Conservation des coefficients 22 et 15 dans la réalisation matérielle de l'électron}

L'exploration historique a d'abord localisé \(p=5\) au moyen d'une comparaison
finie avec \(Z_0\), et a écrit à partir de là
\((27-p,3p)=(22,15)\). Cette procédure conserve sa valeur comme histoire de la
découverte et comme contrôle externe. Ce n'est cependant pas la démonstration
qui gouverne la loi : le dénombrement ultérieur
\[
 \frac{396}{18}=22,
 \qquad
 \frac{270}{18}=15
\]
construit les deux entiers au sein d'APP--TPK et retrouve \(p=5\) comme
conséquence. L'histoire explique comment le couple a été trouvé ; le recensement
explique pourquoi il appartient à l'équation.

\subsection{5ᵉ facteur : le cube de la clôture \(\alpha_{\mathrm{HMT}}\)}

La clôture de \(\alpha_{\mathrm{HMT}}\) précède la mécanique dimensionnelle.
Comme il est établi dans le \cref{chap:alpha}, les douze blocs en base mille
de \(\pi\), \(e\), \(\varphi\) et du sceau \(K\) se combinent au moyen de la
retenue réversible
\[
 P_m+E_m-\Phi_m-K_m+c_{m+1}
 =A_m+1000c_m,
 \qquad c_{13}=c_1=0.
\]
La sortie est
\[
 A_\alpha=
 (007,297,352,569,283,800,997,285,105,472,380,663)
\]
et, par concaténation, la fenêtre dodécaphasique exacte
\[
 \alpha_{12}=\pi_{12}(\alpha_{\mathrm{HMT}})
 =
 \frac{
 7297352569283800997285105472380663
 }{10^{36}}.
\]
Le télescopage des retenues démontre l'identité entière complète ; la
prolongation compatible construit
\(\alpha_{\mathrm{HMT}}=\varprojlim_N A^{(N)}\), et \(\alpha_{12}\) est sa
projection rationnelle à trente-six chiffres. La
comparaison avec la constante de structure fine s'effectue ensuite et
n'intervient pas dans sa construction.

La loi électronique prend la troisième puissance
\[
 \alpha_{\mathrm{HMT}}^3.
\]
Mathématiquement, il s'agit de l'évaluation multiplicative de la troisième
puissance symétrique du même scalaire déjà clos ; on n'introduit ni trois
copies métrologiques ni trois paramètres différents. Sa contribution est
\[
 22\alpha_{\mathrm{HMT}}^3
 =
 0.00000854906599200551727014979807598228\ldots.
\]

\begin{proposition}[Statut du canal cubique électronique]
\label{prop:alpha-cubica-electron-rev6}
Le terme \(22\alpha_{\mathrm{HMT}}^3\) de la loi électronique dépend
uniquement de la clôture préalable de \(\alpha_{\mathrm{HMT}}\), de la
multiplication interne et de l'entier \(22\) du
\cref{thm:coeficientes-electron-rev6}. Il ne dépend pas d'une lecture de la
masse électronique.
\end{proposition}

\begin{proof}
\(\alpha_{12}\) est le rationnel fixé par l'identité de concaténation et de
retenue, et la famille compatible de fenêtres détermine la section
\(\alpha_{\mathrm{HMT}}\). La projection \(\alpha_{12}^3\) certifie les
chiffres imprimés de \(\alpha_{\mathrm{HMT}}^3\).
Le coefficient \(22\) est déterminé par \(396/18\). La composition des
deux opérations ne contient que des données construites antérieurement à la
carte d'énergie.
\end{proof}

Cette proposition fixe exactement ce qui est utilisé dans l'électron. La puissance
trois est le degré cubique de ce canal dans la loi publiée. Il n'est pas nécessaire
de postuler ici une règle universelle selon laquelle toute occurrence
\(\alpha^d\) représente automatiquement une dimension \(d\) ; une
affirmation universelle de ce type exigerait son propre théorème de
multiplicativité entre degrés. La généalogie électronique est complète
sans transformer cette généralisation en prémisse.

\subsection{6ᵉ facteur : torsion minimale et défaut global \(\Delta_4\)}
\index{torsion minimale!\(\Delta_4\)}

Les dénominations \emph{torsion minimale} et \emph{défaut global}
désignent ici un unique invariant, non deux corrections distinctes. L'objet
employé par la loi des masses est
\begin{equation}
 \label{eq:delta4-electron-rev6}
 \boxed{
 \Delta_4
 =
 \frac{\pi-e\log\pi}{270}
 \left(1+\frac{\pi}{729}\right).}
\end{equation}
Ses entrées sont les lectures HMT déjà fixées de \(\pi\) et \(e\), le canal
global \(270\) et la cellule intérieure \(729\). Par conséquent,
\[
 \Delta_4
 =
 0.00011119642269214005405113478612067909\ldots.
\]
Le premier rapport mesure le défaut signé
\(\pi-e\log\pi\) en unités du canal 270 ; le second facteur couple ce
rapport à la cellule intérieure au moyen de \(\pi/729\). Cette lecture n'ajoute pas de
paramètres à la formule : elle sépare ses deux opérations constitutives.
Le \cref{p3-20260826:chap:medida-accion-positiva} montre qu'il s'agit d'une entrée mathématique
prédimensionnelle et la distingue de l'opérateur qui la réalise éventuellement
comme changement de feuillet.

\begin{proposition}[Globalité du défaut électronique]
\label{prop:delta4-global-electron-rev6}
\(\Delta_4\) est un invariant global antérieur à l'espèce. Dans la loi
électronique, il intervient au moyen de
\[
 1+15\Delta_4
 =
 1.00166794634038210081076702179181019\ldots.
\]
Ce n'est pas un correcteur choisi après avoir pris connaissance du résidu de la masse de
l'électron.
\end{proposition}

\begin{proof}
La \cref{eq:delta4-electron-rev6} ne contient ni étiquette d'espèce, ni
masse, ni unité. L'entier \(15\) est déjà construit par \(270/18\).
La composition \(1+15\Delta_4\) est donc déterminée par les données
globales et le recensement interne.
\end{proof}

Il faut distinguer \(\Delta_4\) des autres blocs du caractère de masse et des
autres notions de torsion. En particulier, ce n'est ni \(s_{270}\), ni
l'invariant quadratique \(135\Delta_4^2\), ni le résidu angulaire
\(\varepsilon_{\rm res}\), ni la torsion géométrique de Cartan--Holst.
Substituer l'un à l'autre répéterait ou changerait le niveau de la correction. Dans
l'électron, \(\Delta_4\) conserve la mémoire globale du raffinement ; le facteur
\(15\) déclare combien de fois la famille électronique la lit. Le qualificatif
\emph{minimale} conserve le nom historique de cet invariant et n'est pas utilisé
comme substitut à un théorème variationnel de minimalité.

\subsection{Réalisation matérielle du point fixe (5,5) dans la clôture électronique}

Soit
\[
 \mathcal C_{81}=\{1,\ldots,9\}^2
\]
l'atlas nonadique et
\[
 \rho(r,c)=(10-r,10-c)
\]
son inversion centrale. Le \cref{thm:censo-atlas-81} prouve que l'atlas
contient une unique cellule de génération zéro et que celle-ci porte l'étiquette
positionnelle \(e_{\mathrm{core}}\).

\begin{proposition}[Unicité du centre électronique]
\label{prop:centro-electron-rev6}
L'involution \(\rho\) possède un unique point fixe :
\[
 \boxed{(r,c)=(5,5).}
\]
Son unicité est combinatoire et ne dépend pas de la valeur de
\(\mathfrak e_{\mathrm{HMT}}\).
\end{proposition}

\begin{proof}
L'équation \((r,c)=\rho(r,c)\) équivaut à
\[
 2r=10,\qquad2c=10.
\]
Dans \(\{1,\ldots,9\}^2\), elle a pour unique solution \(r=c=5\). L'évaluation
de la masse n'apparaît pas dans l'argument.
\end{proof}

\begin{corollary}[Universalité de la section centrale]
\label{cor:universalidad-electron-central-rev6}
Tout automorphisme admissible de l'atlas qui entrelace l'inversion centrale
préserve la cellule \((5,5)\).
\end{corollary}

\begin{proof}
Si \(f\rho=\rho f\) et \(\rho(5,5)=(5,5)\), alors
\(\rho f(5,5)=f\rho(5,5)=f(5,5)\). Par conséquent, \(f(5,5)\) est un autre point fixe
de \(\rho\) ; l'unicité impose \(f(5,5)=(5,5)\).
\end{proof}

C'est la raison structurelle pour laquelle la classe électronique est la même dans
toutes les réalisations qui préservent l'atlas et son inversion : un nombre décimal local ne se répète pas
par hasard ; toute carte admissible doit au contraire
transporter le même point fixe et la même chaîne d'invariants. La lecture
physique de cette unicité identifie chaque électron à une réalisation de cette
classe centrale ; la proposition démontre l'universalité interne de l'atlas.

Dans la famille électronique fixée par la construction, ce point est l'ancrage
géométrique de la section. Il convient de ne pas confondre quatre objets :
\[
 (5,5),\qquad
 e_{\mathrm{core}},\qquad
 v_e=(0,0,0,0,0),\qquad
 \text{la signature brute de la cellule centrale}.
\]
Le premier est une position ; le deuxième, son étiquette interne ; le troisième,
l'origine choisie dans le réseau fin d'action ; le quatrième conserve les
données brutes de l'atlas et n'est pas nécessairement le vecteur nul. Leur
coordination explique pourquoi l'électron occupe le centre sans effacer la
généalogie des représentations qui le décrivent.

\paragraph{Routes électroniques orientées \(e_{++},e_{--}\), enveloppe \(e_{\mathrm{env}}\), origine d'action et signatures de la cellule centrale}
La cellule centrale porte deux histoires orientées de douze événements. En écrivant
\(v^2\) pour la concaténation de deux copies d'un sextuplet, la route active
de l'orientation \({++}\), son histoire opposée et l'enveloppe historique
sont respectivement
\[
\begin{aligned}
 e_{\rm act}=e_{++}
   &=(2,2,5,-4,-3,-2)^2,\\
 e_{--}
   &=(4,3,2,-2,-2,-5)^2,\\
 e_{\rm env}
   &=(4,3,5,-4,-3,-5)^2.
\end{aligned}
\tag{E-centro-1}\label{eq:electron-rutas-finas-rev8}
\]
Les deux histoires publient le même mot
\[
 \tau_{12}=\texttt{+++---+++---},
\]
mais ne sont pas la même route. L'enveloppe prend le maximum absolu
coordonnée par coordonnée en conservant le signe ; elle enregistre les amplitudes, mais
aplatit l'ordre généalogique et ne se substitue pas à \(e_{\rm act}\).

Les autres lecteurs du centre vivent dans des codomaines distincts :
\[
\begin{array}{rcll}
 v_e&=&(0,0,0,0,0),
   &\text{origine affine du réseau d’action},\\
 R(5,5)&=&(24,0,12,0,-12),
   &\text{signature APP brute de la cellule},\\
 \operatorname{Sig}^{\Gamma_{108}}_{\rm masa}
   &=&(-12,-4,12,-6,12),
   &\text{projection paire de masse},\\
 \operatorname{Sig}^{\Gamma_{108}}_{\rm carga}
   &=&(0,0,0,0),
   &\text{projection impaire de charge}.
\end{array}
\tag{E-centro-2}\label{eq:electron-firmas-tipadas-rev8}
\]
Le fait que le mot de charge ait une signature impaire nulle ne transforme pas la route en
zéro et n'identifie pas la signature brute à l'origine affine. Position, route
active, route opposée, enveloppe, origine et les deux signatures sont des objets
coordonnés, non six notations pour un même vecteur.

\subsection{Conjugaison particule–antiparticule et persistance de la bande électronique}

L'identité électronique ne s'arrête ni au point central ni à la valeur de la
masse. La conjugaison matérielle se réalise sur des mots dodécaphasiques. Soit
\[
 C_M(\mathbf t)=-\operatorname{rev}(\mathbf t)
\]
l'anti-section orientée et soit
\(\chi\in\{-1,+1\}\) l'orientation de charge. La bande est la classe
\[
 [(\mathbf t,\chi)]_M
 =
 \{(\mathbf t,\chi),(C_M\mathbf t,-\chi)\}
\]
de la \cref{def:banda-dodecafasica-ruta}. La bande conserve ses composantes
paires et présente deux sections de composantes impaires opposées.

\begin{theorem}[Électron et positron comme sections d'une bande]
\label{thm:electron-positron-biografia-rev6}
Pour le caractère de bande du
\cref{thm:principio-conjugacion-masa-banda},
\[
 m(C_M\mathbf t,-\chi)=m(\mathbf t,\chi).
\]
Dans le canal électronique,
\[
 (n,\chi)=(-22,+1)
 \quad\longleftrightarrow\quad
 (+22,-1)
\]
et les deux couples produisent
\[
 \chi n=-22.
\]
L'électron et le positron ont donc la même masse et des orientations de
charge opposées.
\end{theorem}

\begin{proof}
Si \(E_+\) et \(E_-\) sont respectivement les composantes paire et impaire du
caractère par rapport à \(C_M\), alors
\[
 E_+(C_M\mathbf t)=E_+(\mathbf t),
 \qquad
 E_-(C_M\mathbf t)=-E_-(\mathbf t).
\]
D'où,
\[
 E_{\mathrm{band}}(C_M\mathbf t,-\chi)
 =
 E_{\mathrm{band}}(\mathbf t,\chi),
\]
et l'égalité de masse s'obtient en appliquant l'exponentielle. Dans le terme
électronique,
\[
 (+1)(-22)=(-1)(+22)=-22.
\]
\end{proof}

La clôture finie de la \cref{prop:cierre-antiseccion-ct108} renforce cette
interprétation : sur les quatre-vingt-une routes CT--108, ni le signe pur ni
la réversion pure ne préservent la taxonomie, tandis que
\(-\operatorname{rev}\) la préserve bien, avec
\[
 81=54+9+9+9.
\]
L'antiparticule n'est donc ni une masse négative ni une copie
indépendante. C'est l'autre section orientée de la même bande : elle conserve les
invariants pairs qui déterminent la masse et inverse les invariants impairs qui
déterminent charge et orientation.

\subsection{Étape basale de la clôture électronique}

\begin{theorem}[Composition holographique de la section électronique]
\label{thm:cierre-electron-holografico-rev6}
Une fois fixés les objets construits dans les sections précédentes, la section
électronique adimensionnelle est
\[
 \mathfrak e_{\mathrm{HMT}}
 =
 \left\|[P_{\Sigma,3},P_{\Pi,3}]\right\|
 \left[
 \exp\bigl((\mathscr R\circ\mathsf J)(\pi,\varphi)\bigr)
 -\frac{\Sigma(\mathrm{APP})}
 {\Sigma(\mathrm{centro})}\,
 \alpha_{\mathrm{HMT}}^3
 \right]
 \left[
 1+
 \frac{\Sigma(\mathrm{canal})}
 {\Sigma(\mathrm{centro})}\,
 \Delta_4
 \right].
\]
De manière équivalente,
\[
 \boxed{
 \mathfrak e_{\mathrm{HMT}}
 =
 \frac{\sqrt3}{4}
 \left(e^{\varphi/\pi^2}
 -22\alpha_{\mathrm{HMT}}^3\right)
 (1+15\Delta_4).}
\]
Son évaluation est
\[
 \boxed{
 \mathfrak e_{\mathrm{HMT}}
 =
 0.51099892248806607250987087719134943763\ldots.}
\]
Ce scalaire est l'étape basale \(\mathfrak e_e^{(0)}\), non la sortie
électronique directrice. La clôture bidirectionnelle bêta publie
\(0.510998950690000808608\ldots\,\mathrm{MeV}\).
\end{theorem}

\begin{proof}
La \cref{prop:prefactor-app-electron-rev6} apporte
\[
 \left\|[P_{\Sigma,3},P_{\Pi,3}]\right\|=\sqrt3/4.
\]
Le \cref{thm:espejo-electron-rev6} apporte
\[
 (\mathscr R\circ\mathsf J)(\pi,\varphi)=\varphi/\pi^2.
\]
Le \cref{thm:coeficientes-electron-rev6} apporte
\[
 \frac{\Sigma(\mathrm{APP})}{\Sigma(\mathrm{centro})}=22,
 \qquad
 \frac{\Sigma(\mathrm{canal})}{\Sigma(\mathrm{centro})}=15.
\]
La clôture dodécaphasique apporte la fenêtre rationnelle \(\alpha_{12}\) de la
section coinductive \(\alpha_{\mathrm{HMT}}\), et
\cref{eq:delta4-electron-rev6} apporte le défaut global. La substitution
produit la formule abrégée. L'évaluation successive de
\[
\begin{aligned}
 e^{\varphi/\pi^2}
 &=1.17814494259630678081160095680888910\ldots,\\
 22\alpha_{\mathrm{HMT}}^3
 &=0.00000854906599200551727014979807598228\ldots,\\
 1+15\Delta_4
 &=1.00166794634038210081076702179181019\ldots
\end{aligned}
\]
donne la valeur déclarée.
\end{proof}

Le théorème est holographique en un sens vérifiable : chaque facteur local
retient une couche distincte du système complet. Le commutateur retient la
dualité des feuillets ; l'exponentielle retient l'orientation intérieure de la
charnière \(\pi\)--\(\varphi\) ; \(22\) et \(15\) retiennent le recensement global et
son noyau central ; \(\alpha_{\mathrm{HMT}}^3\) retient la clôture fine ; et
\(\Delta_4\) retient le raffinement global. Le centre détermine où
s'ancre la section et la bande détermine comment elle survit sous conjugaison.
La valeur ne remplace pas cette structure : elle est leur évaluation commune.

\subsubsection{Dédoublement de l'action antérieure et postérieure au retour}
\label{sec:electron-desdoblamiento-two-way-rev9}

La réalisation régionale de la coordonnée angulaire conjuguée \(C^*\) distingue deux coordonnées sur la
même droite interne d'action. La coordonnée antérieure au retour est \(H_5\) et
la coordonnée postérieure est
\[
\etaRet
 =H_5-\frac{90}{\pi}\mathscr C_\pi(\alpha_{\rm HMT}).
\]
La notation \(\etaRet\) est la coordonnée sexagésimale canonique définie dans la
branche d'action par \(\etaRet=\Cstar/2-\alphaAngle\). L'égalité précédente est
évaluée dans cette carte ; elle n'introduit pas une seconde coordonnée de retour.
Les cartes additive et multiplicative du dédoublement orienté
sont donc
\begin{equation}
 \boxed{
 \delta_{2w}:=H_5-\etaRet
 =\frac{90}{\pi}\mathscr C_\pi(\alpha_{\rm HMT}),
 \qquad
 R_{\rm act}:=\frac{\hbar^{\rm pre}_{\rm HMT}}
 {\hbar^{\rm ret}_{\rm HMT}}
 =\frac{H_5}{\etaRet}.}
 \label{eq:electron-desdoblamiento-two-way-rev9}
\end{equation}
Les deux coordonnées proviennent d'un unique transport pré-retour/post-retour.
Ce ne sont pas deux constantes de Planck : le facteur commun
\(10^{-34}\mathcal U_S\) fixe la droite dimensionnelle et se simplifie exactement dans
\(R_{\rm act}\). L'ellipse d'action réalise ensuite
\(\operatorname{Area}(\theta)=\hbar_{\rm HMT}^{\rm ret}\theta\), de sorte que
le dédoublement précède la lecture d'holonomie.

Le couplage
\begin{equation}
 \kappa_e^{\beta}
 =80-54\alpha_{\rm HMT}+6\Delta_4,
 \qquad
 \mathfrak e_e^{\beta}
 =\mathfrak e_{\rm HMT}R_{\rm act}^{\kappa_e^{\beta}}
 =0.5109989506900008086\ldots
 \label{eq:electron-beta-two-way-rev9}
\end{equation}
est la clôture quantitative bidirectionnelle de l'électron. Les deux lecteurs
indépendants du registre électronique construits dans le chapitre suivant
sélectionnent, sans consulter la masse de comparaison,
\[
 K_D(\Gamma_e)=K_\Omega(\Gamma_e)=(80,54,6),
\]
et dérivent exactement \(\kappa_e^\beta\). L'entier \(80\) provient de la
première coordonnée de ces lecteurs ; le mot \(w_e=201101\) conserve son
espace de résidence propre dans le canal de la constante d'Euler.

\begin{statusbox}[title={Clôture interne du raffinement bidirectionnel}]
L'identité \eqref{eq:electron-desdoblamiento-two-way-rev9}, la réduction
\((80,54,6)\) et l'équation
\eqref{eq:electron-beta-two-way-rev9} constituent des résultats internes exacts
et sans cible sur le registre électronique. La propagation aux fibres de
rang supérieur se formule au moyen des opérateurs de multisection.
\end{statusbox}

\subsection{Carte énergétique et confrontation de la clôture bidirectionnelle}

Soit \(\mathcal U_E\) une droite réelle de dimension énergétique et soit
\(u_E\) une base choisie. La transduction métrologique est
\[
 \varsigma_{u_E}:\mathbb R\longrightarrow\mathcal U_E,
 \qquad
 \varsigma_{u_E}(x)=x\,u_E.
\]
La carte conserve deux étapes d'une même généalogie. L'étape de base est
\[
 E_e^{(0)}
 =
 \varsigma_{u_E}(\mathfrak e_{\mathrm{HMT}}).
\]
Avec \(u_E=1\,\mathrm{MeV}\),
\[
 E_e^{(0)}
 =
 0.5109989224880660725\ldots\ \mathrm{MeV}.
\]
Le rapport d'action et l'exposant dérivé complètent l'étape de référence,
\[
 E_e^\beta=E_e^{(0)}R_{\rm act}^{\kappa_e^\beta}
 =0.5109989506900008086082571648\ldots\ \mathrm{MeV},
 \qquad
 m_e^\beta=E_e^\beta/c^2.
\]
L'unité n'apparaît pas dans la
\cref{eq:electron-holografico-rev6} ; elle est la base d'une droite dimensionnelle
appliquée à un résultat déjà construit.

La référence CODATA 2022 pour l'énergie équivalente de l'électron est
\cite{CODATA2022}
\[
 E_e^{\mathrm{CODATA}}
 =
 0.51099895069(16)\ \mathrm{MeV}.
\]
Dans cette même carte, le résultat bidirectionnel satisfait
\[
 E_e^\beta-E_e^{\mathrm{CODATA}}
 =
 8.086082571648\times10^{-16}\ \mathrm{MeV},
\]
avec pour différence relative
\[
 1.582406883758\times10^{-15}.
\]
La comparaison est effectuée après fixation de l'évaluation interne. Les
identités du sélecteur de base demeurent
\[
 22=\frac{396}{18},
 \qquad
 15=\frac{270}{18},
 \qquad
 p=27-22=5.
\]
et reçoivent en outre la démonstration combinatoire
\(-22=-|K_e\setminus\iota(R_\pi)|\),
\(+15=|R_\pi\times\mathbb F_3|\). La confrontation métrologique vérifie le
résultat ; le parcours, les lecteurs et le rapport d'action l'engendrent.

\subsection{Conséquences structurelles de la clôture électronique}

La lecture structurelle précédente impose plusieurs distinctions :

\begin{enumerate}
\item \(\sqrt3/4\) est la norme du commutateur APP en dimension trois. Sa
coïncidence avec d'autres formules géométriques ne remplace pas cette provenance.

\item \(\pi/\varphi^2\) et \(\varphi/\pi^2\) ne sont pas interchangeables. Le
premier rapport est l'orientation de clôture ; le second, obtenu au moyen de
\(\mathsf J\), est l'orientation électronique d'autocroissance.

\item \(22\) et \(15\) procèdent de \(396/18\) et de \(270/18\).
Les présenter comme des valeurs logiquement sélectionnées par \(Z_0\) inverserait
la dépendance démontrée.

\item \(\alpha_{\mathrm{HMT}}^3\) utilise la clôture préalable de
\(\alpha_{\mathrm{HMT}}\). La puissance cubique appartient à la loi
électronique ; elle n'autorise pas, à elle seule, une règle universelle
\(\alpha^d\leftrightarrow d\) sans théorème supplémentaire.

\item \(\Delta_4\) est global. Il ne doit pas être remplacé par
\(s_{270}\), par \(135\Delta_4^2\) ou par un résidu particulier de
l'espèce.

\item La position \((5,5)\), l'étiquette \(e_{\mathrm{core}}\), l'origine
fine \(v_e=0\) et la signature brute centrale sont des objets coordonnés, et non
des synonymes.

\item Électron et positron ne se distinguent pas par deux lois de masse : ils sont les
deux sections orientées d'une même bande, de masse paire et de charge impaire.

\item \(0.5109989224\ldots\) est d'abord un scalaire adimensionnel. MeV,
\(\mathrm{MeV}/c^2\), CODATA et \(Z_0\) relèvent de la comparaison
ultérieure.
\end{enumerate}

Ces distinctions expliquent pourquoi l'équation électronique est une
miniature du système HMT--MD. Dans une seule section réapparaissent support,
feuillets, orientation, dénombrement, clôture fine, raffinement, centre, persistance
et conjugaison. Tel est le contenu du mot \emph{holographique} : non
qu'une valeur numérique représente poétiquement le tout, mais que chaque facteur
conserve une dépendance vérifiable à l'égard du tout qui la produit.
